% ===========================================================================
%  Chapitre 16 — Intérêts et escompte
%  PE 2026, composante TRAITEMENT DES DONNÉES — série OSE
% ===========================================================================
\bmtchapitre{Intérêts et escompte}
  {Calculer un intérêt simple, une valeur acquise, un escompte commercial et
   une valeur actuelle, mener un calcul d'intérêt composé, et choisir la
   méthode adaptée à une situation de financement à court ou moyen terme.}
  {Traitement des données \textbullet\ environ 12 heures}

Un capital ne dort jamais tout à fait : placé, il produit un intérêt ;
emprunté, il en coûte un. Les mathématiques financières décrivent ce
mécanisme avec précision, et c'est le langage que parle toute banque, toute
coopérative d'épargne, tout commerçant qui négocie un crédit ou escompte une
traite.

Ce chapitre traite deux grandes familles de calcul. L'\textbf{intérêt
simple}, réservé aux opérations courtes, où l'intérêt ne produit lui-même
aucun intérêt. L'\textbf{intérêt composé}, adapté aux placements plus longs,
où chaque intérêt s'ajoute au capital et se met à son tour à produire des
intérêts. Entre les deux se glisse l'\textbf{escompte}, opération par
laquelle un commerçant obtient aujourd'hui, contre une retenue, l'argent
qu'il n'aurait touché qu'à une échéance future.

% ---------------------------------------------------------------------------
\begin{revision}
Pourcentages et proportionnalité.

\begin{enumerate}[label=\textbf{\arabic*.}, leftmargin=*, itemsep=0.35em]
  \item Calculer $8\,\%$ de $250\,000$.
  \item Un article coûte $60\,000$~Ar après une hausse de $20\,\%$. Quel
        était son prix initial ?
  \item Convertir $8$ mois en douzièmes d'année, puis $90$ jours en
        trois-cent-soixantièmes d'année.
  \item Résoudre $1{,}08^{4} = x$ à l'aide d'une calculatrice ou d'un
        développement.
  \item Résoudre $300 \times (1+t) = 324$ d'inconnue $t$.
\end{enumerate}
\end{revision}

\begin{automatismes}
Sans calculatrice, sans brouillon.

\begin{enumerate}[label=\textbf{\alph*)}, leftmargin=*, itemsep=0.25em]
  \item $10\,\%$ de $500\,000$
  \item $C(1+t)$ pour $C = 100\,000$ et $t = 0{,}05$
  \item Durée de $6$ mois en douzièmes d'année
  \item Durée de $60$ jours en trois-cent-soixantièmes d'année
  \item $1{,}1^{2}$
  \item Valeur acquise d'un capital de $200\,000$ après un intérêt de
        $18\,000$
  \item $(1+t)^{-1}$ pour $t = 0{,}25$
  \item Sens de variation de $C(1+t)^{n}$ en $n$, pour $t>0$
\end{enumerate}
\end{automatismes}

\begin{corrige}[automatismes]
a) $50\,000$ \quad b) $105\,000$ \quad c) $\frac12$ \quad d) $\frac16$
\quad e) $1{,}21$ \quad f) $218\,000$ \quad g) $0{,}8$ \quad h) croissante.
\end{corrige}

% ===========================================================================
\section{Intérêt simple}

\begin{definition}[intérêt simple]
Un capital $C$ placé à \textbf{intérêt simple}, au taux annuel $t$, pendant
une durée $n$ (exprimée en années), produit l'intérêt
\[
  I = C \times t \times n .
\]
La \textbf{valeur acquise} au terme du placement est $C_{n} = C+I
= C(1+tn)$.
\end{definition}

\begin{remarque}
Dans l'intérêt simple, l'intérêt ne porte que sur le capital initial : il ne
se recapitalise jamais. C'est pourquoi ce mode de calcul convient aux
opérations courtes — quelques mois, rarement plus d'un an — comme un crédit
de campagne agricole ou une avance de trésorerie.
\end{remarque}

\begin{methode}{convertir une durée en années}
Une durée exprimée en mois se divise par $12$ ; une durée exprimée en jours
se divise par $360$ dans les usages commerciaux courants (année bancaire de
douze mois de trente jours), sauf précision contraire de l'énoncé.
\end{methode}

\begin{exemple}[intérêt simple et valeur acquise]
Un commerçant emprunte $250\,000$~Ar à intérêt simple, au taux annuel de
$12\,\%$, pendant $8$ mois. La durée vaut $n = \frac{8}{12}$ d'année, donc
\[
  I = 250\,000 \times 0{,}12 \times \frac{8}{12} = 20\,000 \text{ Ar},
\]
et la valeur acquise est $C_{n} = 250\,000+20\,000 = 270\,000$~Ar.
\end{exemple}

\begin{propriete}[retrouver le taux ou la durée]
Les trois grandeurs $C$, $t$, $n$ se déduisent les unes des autres à partir
de $I = Ctn$ :
\[
  t = \frac{I}{Cn}, \qquad n = \frac{I}{Ct}, \qquad C = \frac{I}{tn} .
\]
Le réflexe à avoir est toujours le même : identifier ce que donne l'énoncé,
et isoler l'inconnue dans cette unique formule.
\end{propriete}

\begin{entrainement}[intérêt simple]
\exo[\niveau{1}] Un capital de $180\,000$~Ar est placé à intérêt simple au
taux de $9\,\%$ pendant $5$ mois. Calculer l'intérêt puis la valeur acquise.

\exo[\niveau{1}] Un placement de $6$ mois au taux de $7\,\%$ a rapporté
$10\,500$~Ar d'intérêt. Déterminer le capital placé.

\exo[\niveau{2}] Un capital placé à intérêt simple pendant $9$ mois a permis
d'obtenir une valeur acquise égale à $107\,5\,\%$ du capital initial.
Déterminer le taux annuel.

\exo[\niveau{2}] Deux capitaux, de somme $500\,000$~Ar, sont placés
respectivement $4$ mois au taux de $10\,\%$ et $8$ mois au taux de $8\,\%$ ;
les intérêts produits sont égaux. Déterminer chaque capital.
\end{entrainement}

% ===========================================================================
\section{Escompte commercial}

\begin{definition}[effet de commerce, escompte]
Un \textbf{effet de commerce} (traite, billet à ordre) est un document par
lequel un débiteur s'engage à payer une somme, appelée \textbf{valeur
nominale} $N$, à une date d'\textbf{échéance} fixée. \textbf{Escompter} un
effet, c'est le céder à une banque avant son échéance, contre une somme
immédiate inférieure à $N$ : la différence est l'\textbf{escompte}.
\end{definition}

\begin{theoreme}[escompte commercial]
Pour un effet de valeur nominale $N$, escompté $n$ jours avant l'échéance au
taux annuel $t$, l'\textbf{escompte commercial} et la \textbf{valeur
actuelle commerciale} valent
\[
  e = N \times t \times \frac{n}{360}, \qquad
  a = N-e .
\]
\end{theoreme}

\begin{avertissement}
L'escompte commercial se calcule sur la valeur nominale $N$, et non sur la
somme réellement perçue $a$ : c'est un intérêt simple calculé « à
l'envers », depuis l'échéance vers aujourd'hui. C'est pourquoi le taux
réellement supporté par celui qui escompte est toujours un peu supérieur au
taux annoncé $t$ — un point détaillé dans l'exercice-recherche de ce
chapitre.
\end{avertissement}

\begin{remarque}[escompte rationnel]
Une seconde méthode, l'\textbf{escompte rationnel}, calcule directement la
valeur actuelle par actualisation à intérêt simple :
\[
  a' = \frac{N}{1+t\frac{n}{360}}, \qquad e' = N-a' .
\]
On démontre que $e' \leq e$ : l'escompte rationnel est toujours inférieur ou
égal à l'escompte commercial, à mêmes données. C'est l'escompte commercial
qui est utilisé dans la pratique bancaire courante, et c'est celui que
retient ce manuel sauf mention contraire.
\end{remarque}

\begin{exemple}[escompte commercial]
Un effet de valeur nominale $500\,000$~Ar est escompté $90$ jours avant son
échéance, au taux annuel de $9\,\%$.
\[
  e = 500\,000 \times 0{,}09 \times \frac{90}{360}
  = 500\,000 \times 0{,}0225 = 11\,250 \text{ Ar},
\]
\[
  a = 500\,000-11\,250 = 488\,750 \text{ Ar}.
\]
Le porteur de l'effet reçoit donc $488\,750$~Ar aujourd'hui, contre
$500\,000$~Ar promis dans trois mois.
\end{exemple}

\begin{entrainement}[escompte]
\exo[\niveau{1}] Un effet de $300\,000$~Ar est escompté $60$ jours avant
l'échéance au taux de $10\,\%$. Calculer l'escompte et la valeur actuelle.

\exo[\niveau{2}] Un effet est escompté $120$ jours avant échéance au taux de
$12\,\%$ ; l'escompte vaut $16\,000$~Ar. Déterminer la valeur nominale de
l'effet.

\exo[\niveau{2}] Un commerçant reçoit $291\,000$~Ar pour un effet de valeur
nominale $300\,000$~Ar, escompté au taux de $12\,\%$. Déterminer le nombre de
jours restant avant l'échéance.
\end{entrainement}

% ===========================================================================
\section{Intérêt composé}

\begin{definition}[intérêt composé]
Un capital $C$ placé à \textbf{intérêt composé}, au taux annuel $t$, pendant
$n$ années entières, a pour valeur acquise
\[
  C_{n} = C(1+t)^{n} .
\]
\end{definition}

\begin{remarque}
Ici, chaque intérêt annuel s'incorpore au capital avant de produire à son
tour un intérêt l'année suivante : c'est la « capitalisation ». C'est le mode
de calcul adapté aux placements de plusieurs années — épargne, emprunt
immobilier, obligation — car il reflète correctement la croissance réelle
d'un capital réinvesti.
\end{remarque}

\begin{visualisation}[intérêt simple contre intérêt composé]
\begin{center}
\begin{tikzpicture}[x=1.35cm, y=0.055cm, font=\footnotesize]
  \draw[bmtGris, -{Stealth[length=2.6mm]}] (-0.3,0) -- (6.3,0)
    node[anchor=north west] {$n$ (années)};
  \draw[bmtGris, -{Stealth[length=2.6mm]}] (0,-10) -- (0,340)
    node[anchor=south east] {valeur acquise};
  \draw[bmtVetiver, line width=1.1pt]
    (0,100) -- (1,110) -- (2,120) -- (3,130) -- (4,140) -- (5,150)
    -- (6,160);
  \draw[bmtLaterite, line width=1.1pt]
    (0,100) -- (1,110) -- (2,121) -- (3,133.1) -- (4,146.41)
    -- (5,161.05) -- (6,177.16);
  \node[bmtVetiver, anchor=west] at (6.05,160) {intérêt simple};
  \node[bmtLaterite, anchor=west] at (6.05,177) {intérêt composé};
\end{tikzpicture}
\end{center}

Les deux courbes partent du même capital et du même taux, et coïncident la
première année. Ensuite, l'intérêt simple progresse en ligne droite — chaque
année rapporte exactement le même montant — tandis que l'intérêt composé
s'incurve vers le haut : l'écart entre les deux, faible au début, se creuse
rapidement. C'est pourquoi l'intérêt composé est toujours préféré au-delà
d'un an ou deux.
\end{visualisation}

\begin{propriete}[retrouver le taux ou la durée]
À partir de $C_{n} = C(1+t)^{n}$, en passant au logarithme népérien,
\[
  n = \frac{\ln\left(\dfrac{C_{n}}{C}\right)}{\ln(1+t)},
  \qquad
  t = \left(\frac{C_{n}}{C}\right)^{1/n}-1
    = e^{\frac1n\ln\left(\frac{C_{n}}{C}\right)}-1 .
\]
\end{propriete}

\begin{demonstration}
De $C_{n} = C(1+t)^{n}$ on tire $(1+t)^{n} = \dfrac{C_{n}}{C}$, puis, en
prenant le logarithme népérien de chaque membre,
$n\ln(1+t) = \ln\left(\dfrac{C_{n}}{C}\right)$, d'où la formule de $n$. Pour
le taux, on isole $1+t = \left(\dfrac{C_{n}}{C}\right)^{1/n}
= e^{\frac1n\ln\left(\frac{C_{n}}{C}\right)}$, ce qui donne $t$.
\end{demonstration}

\begin{exemple}[valeur acquise à intérêt composé]
Un capital de $300\,000$~Ar est placé à intérêt composé au taux annuel de
$8\,\%$ pendant $4$ ans.
\[
  C_{4} = 300\,000 \times (1{,}08)^{4}
  = 300\,000 \times 1{,}360\,489 \approx 408\,146{,}69 \text{ Ar}.
\]
\end{exemple}

\begin{remarque}[taux équivalents]
Deux taux, l'un annuel $t$, l'autre relatif à une période plus courte (par
exemple semestrielle $t_{s}$), sont dits \textbf{équivalents} lorsqu'ils
produisent la même valeur acquise sur une année :
$(1+t_{s})^{2} = 1+t$. Ainsi, pour $t = 8\,\%$,
\[
  t_{s} = \sqrt{1{,}08}-1 \approx 0{,}039\,23,
  \qquad \text{soit environ } 3{,}92\,\%.
\]
\end{remarque}

\begin{entrainement}[intérêt composé]
\exo[\niveau{1}] Un capital de $250\,000$~Ar est placé à $6\,\%$ à intérêt
composé pendant $3$ ans. Calculer la valeur acquise.

\exo[\niveau{2}] Un capital placé à $5\,\%$ à intérêt composé double en
combien d'années entières ?

\exo[\niveau{2}] Un capital de $400\,000$~Ar atteint $486\,202{,}82$~Ar après
$4$ ans de placement à intérêt composé. Déterminer le taux annuel.

\exo[\niveau{3}] Déterminer le taux semestriel équivalent à un taux annuel de
$10\,\%$, puis le taux mensuel équivalent à ce même taux annuel.
\end{entrainement}

% ===========================================================================
\begin{annales}
Les mathématiques financières forment l'exercice 3 de l'épreuve, presque
toujours structuré en deux ou trois parties indépendantes : intérêt simple,
puis intérêt composé, puis annuités. L'énoncé qui suit — dont les deux
premières parties relèvent de ce chapitre, la troisième étant traitée au
chapitre suivant — est repris d'un sujet réellement posé.

\exoann{Baccalauréat de l'Enseignement Général, Madagascar, série OSE, option OSE, session 2026 — exercice 3, parties I et II}
\textbf{Partie I.} Un capital $C$ est placé à intérêt simple pendant
$15$ mois ; sa valeur acquise dépasse le capital initial de $10\,\%$.
\begin{enumerate}[label=\textbf{\alph*)}, leftmargin=*, itemsep=0.15em]
  \item Calculer le taux annuel $t$ de ce placement.
  \item Sachant que l'intérêt produit vaut $40\,000$~Ar, déterminer $C$ puis
        la valeur acquise.
\end{enumerate}

\textbf{Partie II.} On place désormais un capital $C = 400\,000$~Ar à
intérêt composé au taux annuel de $6\,\%$, pendant $n$ périodes.
\begin{enumerate}[label=\textbf{\alph*)}, leftmargin=*, itemsep=0.15em]
  \item Déterminer la durée $n$ nécessaire pour que la valeur acquise atteigne
        $535\,290{,}23$~Ar.
  \item En posant $C_{n} = C(1+t)^{n}$, montrer que
        $t = e^{\frac1n\ln\left(\frac{C_{n}}{C}\right)}-1$.
  \item En utilisant cette formule, déterminer le taux annuel nécessaire
        pour que ce même capital atteigne $600\,000$~Ar au bout de $5$~ans.
\end{enumerate}
\end{annales}

\begin{corrige}[annales]
\textbf{Partie I. a)} La valeur acquise vaut $1{,}10\,C$, donc l'intérêt
vaut $0{,}10\,C$. Avec $n = \frac{15}{12} = 1{,}25$ année,
$I = Ctn$ donne $0{,}10\,C = C \times t \times 1{,}25$, d'où
$t = \dfrac{0{,}10}{1{,}25} = 0{,}08$, soit $t = 8\,\%$.

\textbf{b)} $I = 40\,000 = C \times 0{,}08 \times 1{,}25 = 0{,}10\,C$, donc
$C = \dfrac{40\,000}{0{,}10} = 400\,000$~Ar, et la valeur acquise vaut
$400\,000+40\,000 = 440\,000$~Ar.

\textbf{Partie II. a)} On résout $400\,000 \times (1{,}06)^{n}
= 535\,290{,}23$, soit $(1{,}06)^{n} \approx 1{,}338\,226$. Or
$(1{,}06)^{5} = 1{,}338\,226$ : la durée est $n = 5$ périodes (ici, cinq
ans).

\textbf{b)} De $C_{n} = C(1+t)^{n}$ on tire, pour $C \neq 0$,
$(1+t)^{n} = \dfrac{C_{n}}{C}$. En prenant le logarithme népérien, strictement
croissant sur $\Rps$ :
$n\ln(1+t) = \ln\!\left(\dfrac{C_{n}}{C}\right)$, d'où
$\ln(1+t) = \dfrac1n\ln\!\left(\dfrac{C_{n}}{C}\right)$, puis, en composant
par l'exponentielle,
$1+t = e^{\frac1n\ln\left(\frac{C_{n}}{C}\right)}$, ce qui donne bien
$t = e^{\frac1n\ln\left(\frac{C_{n}}{C}\right)}-1$.

\textbf{c)} Avec $C = 400\,000$, $C_{n} = 600\,000$ et $n = 5$ :
\[
  t = e^{\frac15\ln\left(\frac{600\,000}{400\,000}\right)}-1
  = e^{\frac15\ln1{,}5}-1
  = 1{,}5^{1/5}-1 \approx 1{,}084\,47-1 \approx 0{,}0845,
\]
soit un taux annuel voisin de $8{,}45\,\%$.
\end{corrige}

% ===========================================================================
\begin{jecherche}[l'escompte de la traite du riz]
Un collecteur de riz détient une traite de valeur nominale $800\,000$~Ar,
payable dans $120$ jours par un grossiste. Pressé de financer ses achats, il
la fait escompter par sa banque, au taux annuel de $10{,}5\,\%$.

\begin{enumerate}[label=\textbf{\arabic*.}, leftmargin=*, itemsep=0.3em]
  \item Calculer l'escompte commercial, puis la valeur actuelle commerciale
        que la banque doit lui verser.
  \item La banque ajoute en réalité une commission fixe de $5\,000$~Ar,
        prélevée en plus de l'escompte. Calculer les \textbf{agios totaux}
        (escompte $+$ commission) et la somme nette réellement perçue par le
        collecteur.
  \item Montrer que le \textbf{taux réel} supporté par le collecteur — celui
        qui, appliqué seul en intérêt simple sur la somme nette perçue,
        redonnerait la valeur nominale à l'échéance — vaut environ
        $12{,}9\,\%$, sensiblement plus que le taux annoncé de $10{,}5\,\%$.
  \item Calculer, à titre de comparaison, l'escompte rationnel associé aux
        mêmes données, et vérifier qu'il est bien inférieur à l'escompte
        commercial.
  \item Le collecteur pourrait attendre l'échéance sans escompter. Dans
        quelles circonstances économiques ce choix serait-il préférable ?
\end{enumerate}
\end{jecherche}

\begin{corrige}[l'escompte de la traite du riz]
\textbf{1.} $e = 800\,000 \times 0{,}105 \times \dfrac{120}{360}
= 800\,000 \times 0{,}035 = 28\,000$~Ar, et la valeur actuelle commerciale
vaut $a = 800\,000-28\,000 = 772\,000$~Ar.

\textbf{2.} Les agios totaux valent $28\,000+5\,000 = 33\,000$~Ar, et la
somme nette perçue est $800\,000-33\,000 = 767\,000$~Ar.

\textbf{3.} On cherche $t'$ tel que
$767\,000 \times \left(1+t' \times \dfrac{120}{360}\right) = 800\,000$,
soit $1+\dfrac{t'}{3} = \dfrac{800\,000}{767\,000} \approx 1{,}043\,02$,
d'où $\dfrac{t'}{3} \approx 0{,}043\,02$ et $t' \approx 0{,}129$, soit
environ $12{,}9\,\%$. La commission fixe, répartie sur une somme nette plus
faible, alourdit sensiblement le coût réel du crédit par rapport au taux
affiché.

\textbf{4.} L'escompte rationnel vérifie
$a' = \dfrac{800\,000}{1+0{,}105 \times \frac{120}{360}}
= \dfrac{800\,000}{1{,}035} \approx 772\,946$~Ar, donc
$e' = 800\,000-772\,946 \approx 27\,054$~Ar, un peu inférieur à l'escompte
commercial de $28\,000$~Ar — conformément à la propriété du cours.

\textbf{5.} Attendre serait préférable si le collecteur dispose déjà de
liquidités suffisantes pour financer ses achats sans escompter, ou si le taux
d'escompte proposé est nettement supérieur au coût d'un autre financement
disponible : l'escompte n'est avantageux que lorsque le besoin immédiat de
trésorerie dépasse le coût réel du crédit qu'il représente.
\end{corrige}

% ===========================================================================
\begin{jemeteste}
Répondre par vrai ou faux, en justifiant en une ligne.

\begin{enumerate}[label=\textbf{\arabic*.}, leftmargin=*, itemsep=0.28em]
  \item Dans l'intérêt simple, l'intérêt produit lui-même des intérêts.
  \item L'escompte commercial se calcule sur la valeur nominale de l'effet.
  \item L'intérêt composé convient mieux qu'un intérêt simple pour un
        placement de dix ans.
  \item Deux taux équivalents produisent la même valeur acquise sur une
        année.
  \item L'escompte rationnel est toujours supérieur à l'escompte commercial.
  \item $C_{n} = C(1+t)^{n}$ suppose que le taux $t$ reste le même à chaque
        période.
\end{enumerate}
\end{jemeteste}

\begin{corrige}[je me teste]
\textbf{1.} Faux — c'est la définition même de l'intérêt \emph{composé}, pas
simple.
\textbf{2.} Vrai.
\textbf{3.} Vrai — la capitalisation reflète mieux la croissance réelle d'un
capital sur une longue durée.
\textbf{4.} Vrai, par définition.
\textbf{5.} Faux — c'est l'inverse : l'escompte rationnel est toujours
inférieur ou égal à l'escompte commercial.
\textbf{6.} Vrai.
\end{corrige}

% ===========================================================================
\begin{commeaubac}
\bmtsource{Exercice original, au format de l'épreuve — série OSE}
\textbf{Partie I.} Un capital $C$ est placé à intérêt simple pendant
$10$ mois ; sa valeur acquise dépasse le capital initial de $6\,\%$.
\begin{enumerate}[label=\textbf{\alph*)}, leftmargin=*, itemsep=0.15em]
  \item Calculer le taux annuel $t$ de ce placement. \hfill\textup{(0,5 pt)}
  \item Sachant que l'intérêt produit vaut $18\,000$~Ar, déterminer $C$ puis
        la valeur acquise. \hfill\textup{(0,75 pt)}
\end{enumerate}

\textbf{Partie II.} La valeur acquise obtenue en partie I est aussitôt
replacée à intérêt composé, au taux annuel de $5\,\%$.
\begin{enumerate}[label=\textbf{\alph*)}, leftmargin=*, itemsep=0.15em]
  \item Calculer la valeur acquise de ce nouveau placement au bout de
        $4$~ans. \hfill\textup{(0,75 pt)}
  \item Déterminer, en nombre entier d'années, la durée nécessaire pour que
        ce capital double. \hfill\textup{(0,5 pt)}
\end{enumerate}
\end{commeaubac}

\begin{corrige}[comme au baccalauréat]
\textbf{Partie I. a)} La valeur acquise vaut $1{,}06\,C$, donc l'intérêt
vaut $0{,}06\,C$. Avec $n = \frac{10}{12}$ année,
$0{,}06\,C = C \times t \times \frac{10}{12}$, d'où
$t = 0{,}06 \times \frac{12}{10} = 0{,}072$, soit $t = 7{,}2\,\%$.

\textbf{b)} $I = 18\,000 = C \times 0{,}072 \times \frac{10}{12}
= 0{,}06\,C$, donc $C = \dfrac{18\,000}{0{,}06} = 300\,000$~Ar, et la valeur
acquise vaut $300\,000+18\,000 = 318\,000$~Ar.

\textbf{Partie II. a)} $C_{4} = 318\,000 \times (1{,}05)^{4}
= 318\,000 \times 1{,}215\,506 \approx 386\,530{,}99$~Ar.

\textbf{b)} On cherche le plus petit entier $n$ tel que
$(1{,}05)^{n} \geq 2$. Or $(1{,}05)^{14} \approx 1{,}979\,9<2$ et
$(1{,}05)^{15} \approx 2{,}078\,9 \geq 2$ : il faut donc $15$ années pour que
le capital double.
\end{corrige}

% ===========================================================================
\begin{notepourlaclasse}
La difficulté du chapitre n'est jamais le calcul lui-même, mais le choix de
la bonne formule et de la bonne unité de durée. Imposez une phrase de
diagnostic avant tout calcul : simple ou composé ? mois ou jours ? année de
$360$ ou de $365$ jours ? C'est cette lecture, pas l'algèbre, qui distingue
les copies qui réussissent l'exercice 3.

L'écart entre escompte commercial et taux réel supporté — révélé par
l'exercice-recherche — mérite d'être commenté en classe au-delà du calcul :
c'est un vrai point d'économie, qui explique pourquoi les commerçants
négocient leurs commissions bancaires autant que leurs taux.

Enfin, faites toujours vérifier un résultat d'intérêt composé par ordre de
grandeur avant de le rendre définitif : une inversion d'exposant, fréquente
sur $(1+t)^{n}$, produit une erreur souvent visible à l'œil nu si l'élève a
pris l'habitude d'estimer la valeur acquise avant de calculer.
\end{notepourlaclasse}
